\begin{abstract}

Epidemiological models such as SIR provide an interpretable account of how cultural content spreads on streaming platforms, yet they operate on aggregated popularity curves and miss the relational structure among songs, artists, and genres that also shapes that spread. This work asks whether a heterogeneous graph neural network with an explicit temporal component improves the modeling of music diffusion over that population-level baseline. Using the MGD+ data ecosystem combined with daily Spotify chart rankings (Top 200 and Viral 50) for the Brazilian market between 2017 and 2021, we build a heterogeneous graph of 6,526 songs, 1,701 artists, and 530 genres, linked by authorship, genre membership, genre co-occurrence, and popularity co-trajectory relations. On this graph we train \textsc{MusicDiffusionGNN}, which pairs a HeteroGraphSAGE encoder over weekly snapshots with a GRU-based temporal head, and compare it against the SIR baseline under two regimes: retroactive reconstruction of the full curve (Mode~1) and causal $k$-week-ahead forecasting using only past information (Mode~2). The SIR baseline was first reproduced in line with the reference study (viral RMSE $0.0289$, success RMSE $0.0471$, $1,981$ songs), and the constructed graph showed non-random topology, with high clustering and coherent genre communities. In Mode~1, SIR attains lower RMSE because it fits two free parameters to each song's complete curve, a structural advantage of the in-sample regime, though the gap narrows on the weeks a song is actually on the chart. In Mode~2, which matches the intended use of a forecaster, the GNN outperforms SIR across all six evaluated scenarios (two charts $\times$ three horizons), with strong statistical significance (paired Wilcoxon test, $p$ between $6.5\times10^{-4}$ and $5.4\times10^{-32}$) and a lower error for 64\% to 77\% of songs, also surpassing a persistence baseline at longer horizons. These results indicate that the GNN outperforms both baselines at longer horizons, with its margin over SIR widening as the horizon grows; whether this advantage stems specifically from the graph's relational structure, rather than from the added flexibility of a trained neural predictor over a two-parameter fit, remains to be established. Work in progress includes an interpretability analysis of the GNN (edge-type ablation, feature-group importance, and extraction of $\beta$, $\gamma$, and $R_0$ analogs) and a no-graph sequence baseline, to be incorporated in the final version of the dissertation.

\end{abstract}
